% Teil II: der durchgerechnete Weg.

\subsection*{Gegeben und gesucht}
(a)--(c) Terme ausmultiplizieren, (d)--(g) Terme faktorisieren.

\subsection*{a) $(x+1)(x-y+3)$}
\begin{align}
  (x+1)(x-y+3)
  &=x(x-y+3)+1(x-y+3) \\
  &=x^2-xy+3x+x-y+3
  =x^2-xy+4x-y+3 .
\end{align}

\subsection*{b) $(a-x)(x+a)$}
\begin{equation}
  (a-x)(x+a)=(a-x)(a+x)=a^2-x^2 .
\end{equation}

\subsection*{c) $(1-x)(1+x+x^2)$}
\begin{equation}
  (1-x)(1+x+x^2)=1+x+x^2-x-x^2-x^3=1-x^3 .
\end{equation}

\subsection*{d) $x^2+2x+1$}
\begin{equation}
  x^2+2x+1=(x+1)^2 .
\end{equation}

\subsection*{e) $x^2-a^2$}
\begin{equation}
  x^2-a^2=(x-a)(x+a) .
\end{equation}

\subsection*{f) $\tfrac14x^2-3x+9$}
\begin{equation}
  \tfrac14x^2-3x+9=\bigl(\tfrac12x-3\bigr)^2 .
\end{equation}

\subsection*{g) $x^2-8x+15$}
Nullstellen: $x^2-8x+15=0 \Rightarrow x_{1,2}=\frac{8\pm\sqrt{64-60}}{2}
=\frac{8\pm2}{2}$, also $x_1=3$, $x_2=5$:
\begin{equation}
  x^2-8x+15=(x-3)(x-5) .
\end{equation}

\begin{ergebnis}
\begin{align*}
  \text{a)}\ x^2-xy+4x-y+3, &\quad
  \text{b)}\ a^2-x^2, &\quad
  \text{c)}\ 1-x^3,\\
  \text{d)}\ (x+1)^2, &\quad
  \text{e)}\ (x-a)(x+a), &\quad
  \text{f)}\ \bigl(\tfrac12x-3\bigr)^2, &\quad
  \text{g)}\ (x-3)(x-5).
\end{align*}
\end{ergebnis}

\begin{kontrolle}
g) zurückschreiten: $(x-3)(x-5)=x^2-5x-3x+15=x^2-8x+15$; f)
$(\tfrac12x-3)^2=\tfrac14x^2-3x+9$; c) Probe mit $x=2$:
$(1-2)(1+2+4)=-7$, rechts $1-8=-7$. Alles passt.
\end{kontrolle}

\begin{merksatz}
Ausmultiplizieren = Distributivgesetz; Faktorisieren = binomische Formeln
rückwärts oder Nullstellen über die $p$-$q$-Formel bestimmen.
\end{merksatz}
